% META: {"id": "TEXP-MAT-CO-003", "chapitre": "TEXP-MATRICES-MARKOV", "type_objet": "corrige", "exercice_ref": "TEXP-MAT-EX-003", "capacites_codes": ["C3"], "sources_inspiration": [], "status": "generated"}
% BEGIN-VERIFY
% from sympy import Matrix, Rational, eye
% A = Matrix([[Rational(1,2), 0], [Rational(1,4), Rational(3,4)]])
% C = Matrix([[3],[2]])
% I = eye(2)
% X = (I - A).solve(C)
% assert X == Matrix([[6],[14]])
% U0 = Matrix([[0],[0]])
% U1 = A*U0 + C
% U2 = A*U1 + C
% assert U1 == Matrix([[3],[2]])
% assert U2 == Matrix([[Rational(9,2)],[Rational(17,4)]])
% END-VERIFY
\begin{corrige}{TEXP-MAT-EX-003}
\textbf{1.} $U_1=AU_0+C=C=\begin{pmatrix}3\\2\end{pmatrix}$. Puis $U_2=AU_1+C=\begin{pmatrix}\tfrac{1}{2}\times3\\\tfrac{1}{4}\times3+\tfrac{3}{4}\times2\end{pmatrix}+\begin{pmatrix}3\\2\end{pmatrix}=\begin{pmatrix}\tfrac{3}{2}\\\tfrac{9}{4}\end{pmatrix}+\begin{pmatrix}3\\2\end{pmatrix}=\begin{pmatrix}\tfrac{9}{2}\\\tfrac{17}{4}\end{pmatrix}$.

\textbf{2.} $I-A=\begin{pmatrix}\tfrac{1}{2}&0\\-\tfrac{1}{4}&\tfrac{1}{4}\end{pmatrix}$, de déterminant $\tfrac{1}{2}\times\tfrac{1}{4}-0=\tfrac{1}{8}\neq0$ : $I-A$ est inversible. En résolvant $(I-A)X=C$ : $\tfrac{1}{2}x_1=3$ donne $x_1=6$ ; puis $-\tfrac{1}{4}\times6+\tfrac{1}{4}x_2=2$ donne $x_2=14$. Donc $X=\begin{pmatrix}6\\14\end{pmatrix}$.

\textbf{3.} D'après le cours, $U_n=A^n(U_0-X)+X=A^n\begin{pmatrix}-6\\-14\end{pmatrix}+\begin{pmatrix}6\\14\end{pmatrix}$.
\end{corrige}
